% Generated draft: requires mathematical and source review.
\documentclass[11pt]{article}
\usepackage[margin=25mm]{geometry}
\usepackage{fontspec}
\setmainfont{Latin Modern Roman}
\usepackage{amsmath,amssymb,mathtools,unicode-math}
\setmathfont{Latin Modern Math}
\usepackage{enumitem,needspace,xurl,hyperref}
\setlength{\parindent}{0pt}
\setlength{\parskip}{5pt}
\setlength{\emergencystretch}{3em}
\newcommand{\N}{\mathbb{N}}
\newcommand{\Z}{\mathbb{Z}}
\newcommand{\Q}{\mathbb{Q}}
\newcommand{\R}{\mathbb{R}}
\newcommand{\C}{\mathbb{C}}
\newcommand{\F}{\mathbb{F}}
\newcommand{\A}{\mathbb{A}}
\newcommand{\PP}{\mathbb{P}}
\newcommand{\E}{\mathbb{E}}
\newcommand{\eps}{\varepsilon}
\newcommand{\dd}{\,\mathrm{d}}
\begin{document}
\section*{UnsolvedMath ID: 4000005}
\textbf{Source problem code:} catalogue entry
\subsection{Definitions and mathematical statement}
A customary setting for the question is a metric space \((X,d)\) equipped with a Markov kernel \(P\). Write \(P_x=P(x,\cdot)\), assume the relevant measures have finite first moments, and let
\[
 W_1(\mu,\nu)=\inf_{\gamma\in\Pi(\mu,\nu)}\int_{X\times X}d(u,v)\,\gamma(\dd u,\dd v),
 \qquad
 \kappa(x,y)=1-\frac{W_1(P_x,P_y)}{d(x,y)}\quad(x\ne y),
\]
where \(\Pi(\mu,\nu)\) is the set of couplings with the indicated marginals. A uniform positive coarse Ricci lower bound means \(\kappa(x,y)\ge\kappa_0>0\) for every distinct pair in the specified space.

If \(\pi\) is an invariant probability measure, \(\pi P=\pi\), let \(P\) also denote the operator \(Pf(x)=\int f(y)P(x,\dd y)\) on \(L^2(\pi)\), and put
\[
 L^2_0(\pi)=\{f\in L^2(\pi):\textstyle\int f\,\dd\pi=0\},
 \qquad P_0=P|_{L^2_0(\pi)}.
\]
Reversibility means that \(\pi(\dd x)P(x,\dd y)\) is symmetric in \(x,y\). The question asks what quantitative bounds positive coarse Ricci curvature yields when reversibility is absent, including bounds for the spectral radius of \(P_0\), its operator norm \(\|P_0\|_{2\to2}\), or a Poincar\'e inequality. One standard Poincar\'e formulation is
\[
 \operatorname{Var}_{\pi}(f)\le C\,\mathcal E(f,f),\qquad
 \mathcal E(f,f)=\frac12\int_{X\times X}|f(y)-f(x)|^2\,\pi(\dd x)P(x,\dd y),
\]
where \(\operatorname{Var}_{\pi}(f)=\int|f-\int f\,\dd\pi|^2\,\dd\pi\). In a non-reversible setting these spectral-radius, one-step norm, and quadratic-form bounds are different targets and must not be identified automatically. The source does not select a normed function space or a specific Poincar\'e form; \(L^2(\pi)\) and the displayed form are explicit conventional candidates.

A further question is whether finite-state approximation can transfer useful bounds to the general setting. Such a scheme must specify approximating metrics and kernels, preservation of a curvature bound, invariant measures, the mode of convergence, and the passage of the desired estimate to the limit. None is supplied in the excerpt. Likewise the existence of \(\pi\), required regularity, and permissible dependence of constants on parameters other than curvature are not stated there.

The supplied research summary reports partial progress through non-local-gradient Poincar\'e and modified logarithmic Sobolev bounds for non-reversible chains under curvature assumptions, while saying sharp spectral-radius/operator-norm bounds and a general finite-approximation scheme are not fully closed. It supplies no exact inequalities or citations for those results. The 2008 source's open status has no recorded later authoritative resolution check; the later partial-progress label is retained as an attribution.
\subsection{Short English statement}
Determine what positive coarse Ricci curvature controls for non-reversible Markov chains: spectral radius, operator norm, or suitable Poincar\'e inequalities, and whether finite-state approximation transfers such estimates. The precise function space, hypotheses, constants, and approximation scheme remain to be specified from the original source.
\subsection{Sources}
\begin{itemize}
\item Local registry: \url{data/unsolvedmath/problems.json}, permanent UnsolvedMath ID \(4000005\), source code catalogue entry; supplied statement, background, and research summary, updated 2026-08-06.
\item Yann Ollivier, \emph{Discrete Ricci curvature: Open problems} (2008), Problem E: \url{https://www.yann-ollivier.org/rech/publs/problems_curvmarkov.pdf}. Supplied extraction method: PDF text; accessed 2026-07-29. No later authoritative resolution check is recorded.
\item The supplied research summary reports non-local-gradient Poincar\'e and modified logarithmic Sobolev progress without identifiable publications or exact theorem statements.
\end{itemize}
\end{document}
